\chapter{Introduction}\pagestyle{fancy}
ROS-based robotic systems are constructed from loosely coupled nodes that communicate via a publisher–subscriber architecture. Each node typically performs a well-defined task, such as sensor data acquisition, state estimation, or motion planning. Although this separation of concerns promotes modularity, it also requires careful coordination of parameters, topic names, frame hierarchies, and execution dependencies.

Recent developments in cloud robotics, containerization, and orchestration frameworks have improved deployment scalability. Tools such as RobotKube \cite{lampe2023robotkube} and FogROS \cite{fogros22023} facilitate distributed execution and infrastructure management. Meanwhile, language-based systems such as ROScribe \cite{roscribe} demonstrate the feasibility of using large language models (LLMs) to generate code and launch structures from natural language instructions. However, existing approaches either focus primarily on infrastructure orchestration or code generation, and they typically lack a structured reasoning layer capable of translating high-level mission intent into coherent configuration artifacts.

At the same time, LLMs have shown strong capabilities in structured output generation, including JSON schema adherence and code synthesis. This creates an opportunity to explore their integration into robotic configuration pipelines while preserving determinism and structural validity in downstream system assembly.

Despite the modular architecture of ROS, configuring a robotic system remains a highly manual and error-prone process. Small changes in mission objectives—such as switching sensors, enabling mapping, or modifying operational modes—often require updates across multiple YAML files and launch scripts. These modifications must remain internally consistent, particularly with respect to frame identifiers, topic mappings, and dependency relationships between nodes.
Current orchestration tools do not provide a unified mechanism to transform high-level semantic descriptions of robotic tasks into fully consistent ROS deployments. As a result, robotic system configuration 
remains tightly coupled to domain expertise and manual engineering 
effort. The absence of a structured pipeline connecting mission intent
to executable configuration limits reproducibility, slows 
experimentation, and increases the likelihood of integration errors. The problem addressed in this work is how can natural-language mission
descriptions be systematically translated into valid, coherent, and executable
ROS configurations in a deterministic and reproducible manner.